\documentclass[12pt,a4paper,oneside]{article}
\usepackage[utf8]{inputenc}
\usepackage[spanish]{babel}
\usepackage[T1]{fontenc}
\usepackage{lmodern}
\usepackage[margin=2.5cm]{geometry}
\usepackage{amsmath,amssymb}
\usepackage{xcolor}
\usepackage{hyperref}
\usepackage{graphicx}
\usepackage{booktabs}
\usepackage{tabularx}
\usepackage{enumitem}
\usepackage{multirow}
\usepackage{colortbl}
\usepackage{fancyhdr}
\usepackage{lastpage}
\usepackage{tikz}

\hypersetup{colorlinks=true,urlcolor=blue,linkcolor=black}

% Colores corporativos
\definecolor{darkblue}{RGB}{31,56,100}
\definecolor{accent}{RGB}{79,143,247}
\definecolor{success}{RGB}{16,185,129}
\definecolor{warning}{RGB}{245,158,11}
\definecolor{danger}{RGB}{239,68,68}

% Estilos de encabezado
\pagestyle{fancy}
\fancyhf{}
\rhead{\textcolor{accent}{ISO SECURE}}
\lhead{\textcolor{accent}{Etapa 1 - Requerimientos}}
\cfoot{Página \thepage\ de \pageref{LastPage}}
\renewcommand{\headrulewidth}{0.5pt}
\renewcommand{\headrule}{\color{accent}\hrule}

\title{
  {\Huge \textbf{ISO SECURE}}\\
  {\LARGE Plataforma de Cumplimiento ISO 27001}\\[2ex]
  {\Large \textcolor{accent}{Etapa 1: Análisis e Ingeniería de Requerimientos}}\\[1ex]
  {\small Actividad Formativa 1 - 5\%}
}
\author{Agustín Peralta --- Universidad de San Buenaventura}
\date{Entrega: 25 de Mayo de 2026}

\begin{document}

\maketitle

\begin{abstract}
Este documento consolida los requerimientos de seguridad, funcionales y no funcionales del proyecto ISO SECURE, un sistema web para la gestión de seguridad de la información (SGSI) conforme a ISO/IEC 27001:2022. Incluye historias de usuario, casos de uso, casos de abuso y matriz de trazabilidad.
\end{abstract}

\tableofcontents
\newpage

% ============================================================================
\section{Introducción}

ISO SECURE es una plataforma web integral para:
\begin{itemize}
  \item Gestionar y auditar 93 controles ISO 27001 Anexo A
  \item Registrar y analizar incidentes de seguridad
  \item Evaluar riesgos por dominio ISO (matriz Probabilidad x Impacto)
  \item Capacitar a personal en seguridad de la información
  \item Generar reportes ejecutivos y KPIs para la alta dirección
\end{itemize}

\textbf{Stack:} FastAPI (Python 3.13) + React 19 + PostgreSQL 15 (Supabase) + Supabase Auth (JWT)

\textbf{Alcance:} Sistema multi-tenant con 4 roles (admin, auditor, supervisor, analista) y 9 módulos operacionales.

% ============================================================================
\section{Requerimientos de Seguridad}

Los requerimientos de seguridad están alineados con ISO/IEC 27001:2022 y OWASP Top 10 (2024).

\subsection{Confidencialidad (CIA - C)}

\begin{itemize}[label=\textcolor{accent}{[SEC-C]}]
  \item \textbf{SEC-C-1:} Autenticación obligatoria: todo usuario debe autenticarse vía JWT (Supabase OAuth 2.0) antes de acceder a recursos.
  \item \textbf{SEC-C-2:} Control de acceso granular: cada acción (read, write, delete) validada según rol del usuario y pertenencia a empresa (multi-tenant).
  \item \textbf{SEC-C-3:} Cifrado en tránsito: todas las comunicaciones backend-frontend vía HTTPS (TLS 1.3+).
  \item \textbf{SEC-C-4:} Cifrado en reposo: credenciales, tokens sensibles cifrados en BD (Supabase built-in).
  \item \textbf{SEC-C-5:} CORS restringido: solo dominios autorizados pueden hacer requests (whitelist en FastAPI).
\end{itemize}

\subsection{Integridad (CIA - I)}

\begin{itemize}[label=\textcolor{success}{[SEC-I]}]
  \item \textbf{SEC-I-1:} Validación de entrada: todos los payloads validados con Pydantic v2 (type hints + constraints).
  \item \textbf{SEC-I-2:} Prevención SQL injection: SQLAlchemy ORM con prepared statements (no raw SQL).
  \item \textbf{SEC-I-3:} CSRF protection: tokens CSRF en formularios críticos (si aplica SSR; React con Axios maneja).
  \item \textbf{SEC-I-4:} Integridad de datos: checksums o versionado para cambios en controles y riesgos.
  \item \textbf{SEC-I-5:} Auditoría de cambios: toda modificación (create, update, delete) registrada en tabla de auditoría con timestamp, usuario, acción y deltas.
\end{itemize}

\subsection{Disponibilidad (CIA - A)}

\begin{itemize}[label=\textcolor{danger}{[SEC-A]}]
  \item \textbf{SEC-A-1:} Rate limiting: máximo 100 requests/minuto por IP; 50/min por usuario autenticado.
  \item \textbf{SEC-A-2:} Timeout de sesión: JWT expira en 1 hora; refresh token en 7 días.
  \item \textbf{SEC-A-3:} Resiliencia: base de datos replicada (Supabase built-in); backend sin estado (stateless).
  \item \textbf{SEC-A-4:} Logging y monitoreo: registro centralizado de eventos de seguridad (intentos de acceso fallidos, cambios de permisos).
\end{itemize}

\subsection{Adicionales (OWASP)}

\begin{itemize}[label=\textcolor{warning}{[SEC-X]}]
  \item \textbf{SEC-X-1:} Content Security Policy (CSP): header restrictivo (`default-src 'self'`, etc.) contra XSS.
  \item \textbf{SEC-X-2:} X-Frame-Options: `DENY` para prevenir clickjacking.
  \item \textbf{SEC-X-3:} X-Content-Type-Options: `nosniff` para prevenir MIME sniffing.
  \item \textbf{SEC-X-4:} HSTS: Strict-Transport-Security con max-age=31536000.
  \item \textbf{SEC-X-5:} Permissions-Policy: desabilitar cámara, micrófono, geolocalización.
\end{itemize}

% ============================================================================
\section{Requerimientos Funcionales}

\subsection{RF-1: Autenticación y Gestión de Usuarios}

\begin{tabularx}{\linewidth}{|l|X|}
\hline
\textbf{ID} & RF-1.1 \\
\textbf{Nombre} & Registro de usuario \\
\textbf{Descripción} & Un nuevo usuario puede registrarse con email y contraseña; se valida unicidad de email y se envía token de confirmación. \\
\textbf{Actores} & Usuario anónimo \\
\textbf{Precondición} & Email no registrado previamente \\
\textbf{Postcondición} & Perfil de usuario creado; email de confirmación enviado \\
\hline
\end{tabularx}

\begin{tabularx}{\linewidth}{|l|X|}
\hline
\textbf{ID} & RF-1.2 \\
\textbf{Nombre} & Login \\
\textbf{Descripción} & Usuario autenticado recibe JWT token (válido 1 hora) y refresh token (válido 7 días). \\
\textbf{Actores} & Usuario registrado \\
\textbf{Precondición} & Email y contraseña correctos \\
\textbf{Postcondición} & Tokens almacenados en localStorage; usuario redirigido a dashboard \\
\hline
\end{tabularx}

\begin{tabularx}{\linewidth}{|l|X|}
\hline
\textbf{ID} & RF-1.3 \\
\textbf{Nombre} & Gestión de roles y permisos \\
\textbf{Descripción} & Admin puede asignar roles a usuarios (admin, auditor, supervisor, analista). Cada rol tiene permisos sobre 11 vistas. \\
\textbf{Actores} & Admin \\
\textbf{Precondición} & Usuario existente; admin autenticado \\
\textbf{Postcondición} & Rol asignado; permisos actualizados dinámicamente \\
\hline
\end{tabularx}

\subsection{RF-2: Gestión de Controles ISO 27001}

\begin{tabularx}{\linewidth}{|l|X|}
\hline
\textbf{ID} & RF-2.1 \\
\textbf{Nombre} & Consultar catálogo de controles \\
\textbf{Descripción} & Auditor/Supervisor visualiza 93 controles ISO Anexo A filtrados por empresa, dominio (A.5 a A.18) y estado (no_iniciado, en_proceso, implementado, etc.). \\
\textbf{Actores} & Auditor, Supervisor \\
\textbf{Precondición} & Usuario autenticado \\
\textbf{Postcondición} & Lista de controles mostrada con filtros y búsqueda \\
\hline
\end{tabularx}

\begin{tabularx}{\linewidth}{|l|X|}
\hline
\textbf{ID} & RF-2.2 \\
\textbf{Nombre} & Actualizar estado de control \\
\textbf{Descripción} & Supervisor puede cambiar el estado de un control (no_iniciado -> en_proceso -> implementado -> auditado). \\
\textbf{Actores} & Supervisor \\
\textbf{Precondición} & Control existente; Supervisor autenticado \\
\textbf{Postcondición} & Estado actualizado; auditoría registrada \\
\hline
\end{tabularx}

\begin{tabularx}{\linewidth}{|l|X|}
\hline
\textbf{ID} & RF-2.3 \\
\textbf{Nombre} & Calcular indicador de cumplimiento \\
\textbf{Descripción} & Sistema calcula automáticamente el porcentaje de cumplimiento como (controles\_implementados / total\_controles) x 100. \\
\textbf{Actores} & Sistema \\
\textbf{Precondición} & Al menos 1 control registrado \\
\textbf{Postcondición} & KPI mostrado en dashboard \\
\hline
\end{tabularx}

\subsection{RF-3: Gestión de Incidentes de Seguridad}

\begin{tabularx}{\linewidth}{|l|X|}
\hline
\textbf{ID} & RF-3.1 \\
\textbf{Nombre} & Registrar incidente \\
\textbf{Descripción} & Analista crea incidente con: descripción, severidad (1-5), empresa, fecha de detección, estado inicial (abierto). \\
\textbf{Actores} & Analista \\
\textbf{Precondición} & Incidente detectado \\
\textbf{Postcondición} & Incidente almacenado; notificación a auditor \\
\hline
\end{tabularx}

\begin{tabularx}{\linewidth}{|l|X|}
\hline
\textbf{ID} & RF-3.2 \\
\textbf{Nombre} & Calcular MTTR (Mean Time To Resolve) \\
\textbf{Descripción} & Para incidentes cerrados, MTTR = fecha\_cierre - fecha\_apertura (en horas). Dashboard muestra promedio por empresa. \\
\textbf{Actores} & Sistema \\
\textbf{Precondición} & Incidente cerrado \\
\textbf{Postcondición} & KPI MTTR actualizado \\
\hline
\end{tabularx}

\subsection{RF-4: Evaluación de Riesgos}

\begin{tabularx}{\linewidth}{|l|X|}
\hline
\textbf{ID} & RF-4.1 \\
\textbf{Nombre} & Crear evaluación de riesgo \\
\textbf{Descripción} & Auditor registra riesgo con: descripción, probabilidad (1-5), impacto (1-5), dominio ISO, estado. Riesgo = P x I (1-25). \\
\textbf{Actores} & Auditor \\
\textbf{Precondición} & Riesgo identificado \\
\textbf{Postcondición} & Riesgo registrado; matriz de riesgo actualizada \\
\hline
\end{tabularx}

\begin{tabularx}{\linewidth}{|l|X|}
\hline
\textbf{ID} & RF-4.2 \\
\textbf{Nombre} & Visualizar matriz de riesgo \\
\textbf{Descripción} & Dashboard muestra matriz 5x5 (PxI) con color: verde (1-5), amarillo (6-12), rojo (13-25). Drill-down a riesgos por celda. \\
\textbf{Actores} & Admin, Auditor \\
\textbf{Precondición} & Al menos 1 riesgo registrado \\
\textbf{Postcondición} & Matriz visualizada \\
\hline
\end{tabularx}

\subsection{RF-5: Capacitaciones}

\begin{tabularx}{\linewidth}{|l|X|}
\hline
\textbf{ID} & RF-5.1 \\
\textbf{Nombre} & Acceder a cursos de seguridad \\
\textbf{Descripción} & Usuario puede ver catálogo de cursos (SGSI 101, ISO 27001 Basics, Phishing Awareness, etc.), reproducir video y marcar como completado. \\
\textbf{Actores} & Cualquier usuario autenticado \\
\textbf{Precondición} & Cursos cargados en BD \\
\textbf{Postcondición} & Progreso registrado \\
\hline
\end{tabularx}

\subsection{RF-6: Auditoría Interna}

\begin{tabularx}{\linewidth}{|l|X|}
\hline
\textbf{ID} & RF-6.1 \\
\textbf{Nombre} & Ejecutar checklist de auditoría \\
\textbf{Descripción} & Auditor marca items de checklist pre-configurado (ej: ``¿Se implementó control A.5.1?''). Sistema calcula % cumplimiento. \\
\textbf{Actores} & Auditor \\
\textbf{Precondición} & Checklist definido \\
\textbf{Postcondición} & Resultados almacenados; informe generado \\
\hline
\end{tabularx}

\subsection{RF-7: Reportes y Exportación}

\begin{tabularx}{\linewidth}{|l|X|}
\hline
\textbf{ID} & RF-7.1 \\
\textbf{Nombre} & Exportar reporte ejecutivo \\
\textbf{Descripción} & Admin/Auditor genera PDF con: KPIs, matriz de riesgos, tabla de controles, gráficos MTTR, progreso capacitación. \\
\textbf{Actores} & Admin, Auditor \\
\textbf{Precondición} & Datos disponibles \\
\textbf{Postcondición} & PDF descargado; webhook a n8n para procesamiento \\
\hline
\end{tabularx}

% ============================================================================
\section{Requerimientos No Funcionales}

\begin{tabularx}{\linewidth}{|l|X|}
\hline
\textbf{RNF-1} & \textbf{Performance} --- Latencia máxima 200ms (p95) para GET; 500ms (p95) para POST. \\
\hline
\textbf{RNF-2} & \textbf{Escalabilidad} --- Capacidad para 1000 usuarios concurrentes por empresa; 100 empresas máximo (MVP). \\
\hline
\textbf{RNF-3} & \textbf{Disponibilidad} --- Uptime 99.5\% (SLA Supabase). Recovery Time Objective (RTO) \textless{} 1 hora. \\
\hline
\textbf{RNF-4} & \textbf{Usabilidad} --- Interfaz responsive (mobile, tablet, desktop); accesibilidad WCAG 2.1 AA mínimo. \\
\hline
\textbf{RNF-5} & \textbf{Mantenibilidad} --- Código documentado (docstrings); commits atómicos; CI/CD con tests automáticos. \\
\hline
\textbf{RNF-6} & \textbf{Portabilidad} --- Docker containerizado; deployment multi-cloud compatible. \\
\hline
\textbf{RNF-7} & \textbf{Confiabilidad} --- MTBF \textgreater{} 720 horas; MTTR \textless{} 2 horas para incident crítico. \\
\hline
\textbf{RNF-8} & \textbf{Compatibilidad} & Chrome 120+, Firefox 121+, Safari 17+, Edge 120+. \\
\hline
\end{tabularx}

% ============================================================================
\section{Historias de Usuario}

\subsection{HU-1: Auditor - Consultar cumplimiento por empresa}

\begin{description}
  \item[Como] auditor de seguridad
  \item[Quiero] visualizar el porcentaje de cumplimiento de controles ISO 27001 para cada empresa en el sistema
  \item[Para] identificar rápidamente cuál organización necesita más esfuerzo en implementación
\end{description}

\textbf{Criterios de Aceptación:}
\begin{enumerate}
  \item Mostrar tabla con: Empresa | Total Controles | Implementados | \% Cumplimiento | Tendencia
  \item Filtro por rango de fechas
  \item Exportar a CSV
  \item Drill-down: clic en empresa abre detalle de controles por estado
\end{enumerate}

\subsection{HU-2: Supervisor - Cambiar estado de control}

\begin{description}
  \item[Como] supervisor de implementación
  \item[Quiero] actualizar el estado de un control ISO (no\_iniciado -> en\_proceso -> implementado -> auditado)
  \item[Para] llevar un seguimiento de progreso y comunicar avances al team
\end{description}

\textbf{Criterios de Aceptación:}
\begin{enumerate}
  \item Dropdown con transiciones válidas (no todos los estados son alcanzables de uno)
  \item Comentario obligatorio al cambiar de estado
  \item Auditoría: registrar quién, cuándo, de qué a qué
  \item Notificación a auditor si hay cambio crítico
\end{enumerate}

\subsection{HU-3: Analista - Registrar incidente de seguridad}

\begin{description}
  \item[Como] analista de incidentes
  \item[Quiero] registrar un incidente detectado (malware, acceso no autorizado, leak, etc.)
  \item[Para] mantener historial de eventos y calcular métricas (MTTR, tasa de resolución)
\end{description}

\textbf{Criterios de Aceptación:}
\begin{enumerate}
  \item Formulario con: descripción, severidad (1-5), fecha detección, empresa afectada, dominio ISO relacionado (opcional)
  \item Validación: severidad requerida; descripción mín 10 caracteres
  \item Email de notificación al auditor
  \item Link a reporte de incidente en 24h
\end{enumerate}

\subsection{HU-4: Admin - Gestionar usuarios y roles}

\begin{description}
  \item[Como] administrador del sistema
  \item[Quiero] crear, editar, eliminar usuarios y asignarles roles
  \item[Para] controlar quién accede a qué datos en cada empresa
\end{description}

\textbf{Criterios de Aceptación:}
\begin{enumerate}
  \item CRUD completo: crear (email validado), editar (nombre, rol, empresa), desactivar (no borrar)
  \item Visualización de último login
  \item Bulk actions: asignar rol a múltiples usuarios
  \item Auditoría completa de cambios de rol
\end{enumerate}

\subsection{HU-5: Usuario - Realizar capacitación}

\begin{description}
  \item[Como] empleado de la organización
  \item[Quiero] acceder a cursos de seguridad (SGSI 101, Phishing Awareness, etc.)
  \item[Para] cumplir con el requerimiento de capacitación obligatoria
\end{description}

\textbf{Criterios de Aceptación:}
\begin{enumerate}
  \item Catálogo de cursos con video embebido (YouTube o similar)
  \item Progreso visual (barra \%)
  \item Certificado descargable al completar
  \item Recordatorio si no completa en 30 días
\end{enumerate}

\subsection{HU-6: Auditor - Evaluar riesgo}

\begin{description}
  \item[Como] auditor de riesgos
  \item[Quiero] registrar una evaluación de riesgo (descripción, probabilidad, impacto, dominio ISO)
  \item[Para] alimentar la matriz de riesgos y generar hoja de ruta de mitigación
\end{description}

\textbf{Criterios de Aceptación:}
\begin{enumerate}
  \item Formulario: descripción, P (1-5), I (1-5), dominio ISO, estado
  \item Cálculo automático de riesgo (P x I)
  \item Matriz 5x5 actualizada en tiempo real
  \item Sugerencia de mitigaciones basadas en riesgo
\end{enumerate}

\subsection{HU-7: Admin --- Generar reporte ejecutivo}

\begin{description}
  \item[Como] director de seguridad
  \item[Quiero] exportar un reporte PDF con KPIs, gráficos, matriz de riesgos para la junta directiva
  \item[Para] documentar el estado del SGSI y justificar inversión en seguridad
\end{description}

\textbf{Criterios de Aceptación:}
\begin{enumerate}
  \item PDF con: portada, indice, KPIs (cumplimiento \%, MTTR, tasa incidentes), gráficos Chart.js, tabla de riesgos, tabla de controles críticos
  \item Watermark ``Confidencial''
  \item Timestamp de generación
  \item Opción de programar generación mensual
\end{enumerate}

% ============================================================================
\section{Casos de Uso}

\subsection{CU-1: Implementar un control ISO}

\textbf{Actores:} Supervisor (principal), Auditor (secundario), Sistema

\textbf{Precondición:} Control existe; Supervisor autenticado

\textbf{Flujo Principal:}
\begin{enumerate}
  \item Supervisor abre módulo ``Controles''
  \item Selecciona control con estado ``no\_iniciado''
  \item Hace clic en ``Cambiar Estado'' -> ``En Proceso''
  \item Comenta: ``Iniciado implementación de A.5.1 Control de Acceso''
  \item Sistema registra cambio; notifica a Auditor
  \item Supervisor cambia a ``Implementado'' cuando finaliza
  \item Auditor revisa; cambia a ``Auditado'' si aprueba
\end{enumerate}

\textbf{Flujo Alternativo (A1):} Supervisor añade evidencia (documento, screenshot)
\begin{enumerate}
  \item En paso 3, hace upload de archivo (PDF, JPG máx 5MB)
  \item Sistema valida y almacena; crea thumbnail para preview
\end{enumerate}

\textbf{Postcondición:} Control en estado ``Auditado''; auditoría registrada; KPI cumplimiento actualizado

\subsection{CU-2: Gestionar incidente de seguridad}

\textbf{Actores:} Analista (principal), Auditor (secundario), Sistema, Correo

\textbf{Precondición:} Incidente detectado; Analista autenticado

\textbf{Flujo Principal:}
\begin{enumerate}
  \item Analista accede a ``Incidentes'' -> ``Crear Nuevo''
  \item Rellena: título, descripción, severidad (ej: 4), empresa afectada, fecha detección
  \item Selecciona dominio ISO relacionado (A.12.3 Logging)
  \item Envía; sistema crea ID único (ej: INC-2026-001)
  \item Auditor recibe notificación por email
  \item Analista/Auditor registran acciones de mitigación en pestaña ``Actividades''
  \item Cuando resuelven, cambian a ``Cerrado''; sistema calcula MTTR automáticamente
\end{enumerate}

\textbf{Flujo Alternativo (A1):} Incidente escala a crítico
\begin{enumerate}
  \item Si severidad \textgreater{}=\textgreater{} 4, sistema envía SMS y llamada a on-call
\end{enumerate}

\textbf{Postcondición:} Incidente registrado; MTTR calculado; KPI de incidentes actualizado

\subsection{CU-3: Realizar auditoría interna}

\textbf{Actores:} Auditor (principal), Sistema

\textbf{Precondición:} Checklist definido; Auditor autenticado

\textbf{Flujo Principal:}
\begin{enumerate}
  \item Auditor abre ``Auditoría'' -> ``Nueva Auditoría Interna''
  \item Selecciona empresa y período (ej: Abril 2026)
  \item Abre checklist pre-cargado (15 items)
  \item Para cada item, marca: Cumple / No Cumple / Parcial / No Aplica
  \item Agrega comentarios y evidencia (links)
  \item Calcula automáticamente \% cumplimiento
  \item Genera borrador de reporte
  \item Envía para aprobación del Auditor Jefe
  \item Una vez aprobado, marca como ``Finalizado''
\end{enumerate}

\textbf{Postcondición:} Reporte de auditoría registrado; KPI de auditoría actualizado; notificación a Auditor Jefe

% ============================================================================
\section{Casos de Abuso}

\subsection{CA-1: Atacante intenta SQL injection}

\textbf{Actor Malicioso:} Atacante externo

\textbf{Contexto:} API REST en \texttt{/api/controls}

\textbf{Escenario de Ataque:}
\begin{enumerate}
  \item Atacante intenta: \texttt{GET /api/controls?filter=id=1 OR 1=1}
  \item Sistema debe rechazar con error 400 (Bad Request)
\end{enumerate}

\textbf{Mitigación:}
\begin{itemize}
  \item [OK] SQLAlchemy ORM con prepared statements (no raw SQL)
  \item [OK] Validación Pydantic en todos los inputs
\end{itemize}

\subsection{CA-2: Usuario intenta acceder a datos de otra empresa}

\textbf{Actor Malicioso:} Usuario autenticado (insider)

\textbf{Contexto:} Multi-tenancy

\textbf{Escenario de Ataque:}
\begin{enumerate}
  \item Usuario de Empresa A intenta: \texttt{GET /api/incidents?empresa\_id=999} (Empresa B)
  \item Sistema debe rechazar con 403 (Forbidden)
\end{enumerate}

\textbf{Mitigación:}
\begin{itemize}
  \item [OK] Middleware de autorización: valida empresa\_id del token vs. parámetro
  \item [OK] Row-level security en BD (Supabase RLS)
\end{itemize}

\subsection{CA-3: Admin deshabilitado intenta reasignarse rol}

\textbf{Actor Malicioso:} Admin cuya cuenta fue desactivada

\textbf{Contexto:} JWT válido pero usuario desactivado

\textbf{Escenario de Ataque:}
\begin{enumerate}
  \item Admin desactivado recibe JWT antes de ser desactivado
  \item Intenta: \texttt{POST /api/users/me/role} con rol=admin
  \item Sistema debe rechazar verificando campo activo=False
\end{enumerate}

\textbf{Mitigación:}
\begin{itemize}
  \item [OK] En cada request, verificar estado de usuario (en caché o BD)
  \item [OK] Invalidar JWT si usuario está inactivo
\end{itemize}

\subsection{CA-4: Auditor intenta modificar incidente ajeno}

\textbf{Actor Malicioso:} Auditor de Empresa A

\textbf{Contexto:} RBAC insuficiente

\textbf{Escenario de Ataque:}
\begin{enumerate}
  \item Auditor A intenta: \texttt{PATCH /api/incidents/123} (pertenece a Empresa B)
  \item Sistema debe rechazar con 403
\end{enumerate}

\textbf{Mitigación:}
\begin{itemize}
  \item [OK] Verificar owner (empresa\_id) en middleware
  \item [OK] Auditor solo puede leer/modificar datos de su propia empresa
\end{itemize}

\subsection{CA-5: Token JWT interceptado (MITM)}

\textbf{Actor Malicioso:} Atacante de red (MITM)

\textbf{Contexto:} Comunicación cliente-servidor

\textbf{Escenario de Ataque:}
\begin{enumerate}
  \item Atacante intercepta JWT en tránsito
  \item Intenta reutilizarlo desde otra IP
  \item Sistema debe detectar anomalía o expiración
\end{enumerate}

\textbf{Mitigación:}
\begin{itemize}
  \item [OK] HTTPS obligatorio (TLS 1.3)
  \item [OK] JWT con expiry corto (1 hora)
  \item [OK] Refresh token separado (7 días, HttpOnly)
  \item [OK] IP binding opcional (Supabase setting)
\end{itemize}

\subsection{CA-6: XSS en campo de comentarios}

\textbf{Actor Malicioso:} Supervisor

\textbf{Contexto:} Formulario de ``Comentario en Control''

\textbf{Escenario de Ataque:}
\begin{enumerate}
  \item Supervisor ingresa: \texttt{\textless{}img src=x onerror="alert('XSS')"\textgreater{}}
  \item Si no se sanitiza, ejecuta JavaScript en navegador de Auditor
\end{enumerate}

\textbf{Mitigación:}
\begin{itemize}
  \item [OK] Sanitización en backend (Bleach o similar) antes de guardar
  \item [OK] CSP header: \texttt{script-src 'self'} (bloquea inline scripts)
  \item [OK] React escapa valores por defecto en JSX
\end{itemize}

% ============================================================================
\section{Matriz de Trazabilidad}

La siguiente tabla relaciona Historias de Usuario (HU) con Requerimientos Funcionales (RF) y Controles ISO.

\begin{tiny}
\begin{tabularx}{\textwidth}{|l|l|l|l|}
\hline
\textbf{HU} & \textbf{RF Relacionados} & \textbf{Controles ISO} & \textbf{Endpoints} \\
\hline
HU-1 & RF-2.1, RF-2.3 & A.5.15, A.9.2 & GET /api/controls, GET /dashboard/kpis \\
\hline
HU-2 & RF-2.2, RF-6.1 & A.5.15, A.9.4 & PATCH /api/controls/\{id\}, POST /api/auditoria \\
\hline
HU-3 & RF-3.1, RF-3.2 & A.12.4, A.16.1 & POST /api/incidents, GET /api/incidents \\
\hline
HU-4 & RF-1.3 & A.5.1, A.6.2 & POST/PATCH/DELETE /api/users, GET /api/users \\
\hline
HU-5 & RF-5.1 & A.7.2, A.7.3 & GET /api/capacitaciones, PATCH /api/curso-progreso \\
\hline
HU-6 & RF-4.1, RF-4.2 & A.12.6, A.14.1 & POST /api/risk, GET /api/risk/matriz \\
\hline
HU-7 & RF-7.1 & A.5.1, A.16.1 & POST /api/reportes/export-pdf \\
\hline
\end{tabularx}
\end{tiny}

\newpage

\section{Métricas de Cobertura}

\begin{itemize}
  \item \textbf{Historias de Usuario:} 7 HU definidas
  \item \textbf{Casos de Uso:} 3 CU principales
  \item \textbf{Casos de Abuso:} 6 CA identificados
  \item \textbf{Requerimientos Funcionales:} 14 RF (7 módulos)
  \item \textbf{Requerimientos No Funcionales:} 8 RNF
  \item \textbf{Requerimientos de Seguridad:} 16 SEC (CIA + OWASP)
  \item \textbf{Endpoints API:} 39 endpoints REST
  \item \textbf{Tablas BD:} 9 tablas (multi-tenant)
  \item \textbf{Controles ISO:} 93 controles (Anexo A)
\end{itemize}

\section{Conclusiones}

Los requerimientos de ISO SECURE cubren un espectro completo desde seguridad (17 controles) hasta funcionalidad (14 requerimientos), con énfasis en:
\begin{enumerate}
  \item \textbf{Seguridad by Design:} CIA + OWASP integrados
  \item \textbf{Multi-tenancy:} Aislamiento de datos por empresa
  \item \textbf{RBAC granular:} 4 roles con 11 vistas diferenciadas
  \item \textbf{Trazabilidad completa:} HU <-> RF <-> Controles ISO <-> Endpoints
  \item \textbf{Resiliencia:} Auditoría de cambios, MTTR/MTBF, RLS en BD
\end{enumerate}

Este documento servirá como referencia para las Etapas 2-5 (Arquitectura, Codificación, Pruebas, Entrega).

\vfill

\begin{center}
  \textit{Documento generado: \today}\\
  \textit{Etapa 1 - Análisis e Ingeniería de Requerimientos}\\
  \textit{Universidad de San Buenaventura - Ingeniería de Sistemas}
\end{center}

\end{document}
